% 定理 thm:multiclass-growth-bound 证明中的前缀分组，在六模式族 V 上复算。
% m=2，末坐标为 x2；前缀 u 为 x1 上的标签（网格的行）。
% 行“正常”可接{正常,广告}，行“广告”可接{广告,诈骗}，行“诈骗”可接{正常,诈骗}；
% r(u)=2，故 |V|=Σr(u)=6，|U|=3，每个 U_{a,b} 恰含一个前缀，
% |V| ≤ |U| + Σ|U_{a,b}| = 3+3 = 6（此例取等号）。每个 U_{a,b} 只有一个前缀，Natarajan 维 0 = d-1。
\begingroup
\begin{tikzpicture}[figure picture,x=1mm,y=1mm]
 \def\cs{10}
 \foreach \r in {1,2,3}{\foreach \c in {1,2,3}{
   \draw[FigureGrid,line width=.5pt] ({10+\cs*(\c-1)},{20-\cs*(\r-1)}) rectangle ++(\cs,\cs);}}
 \foreach \lab [count=\c] in {正常,广告,诈骗}{
   \node[anchor=south,font=\figurefont\scriptsize] at ({10+\cs*(\c-.5)},30.5) {\lab};}
 \foreach \lab [count=\r] in {正常,广告,诈骗}{
   \node[anchor=east,font=\figurefont\scriptsize] at (9,{25-\cs*(\r-1)}) {\lab};}
 \node[figure note,anchor=south] at (25,35) {末坐标 $x_2$};
 \node[figure note,anchor=east] at (9,35) {前缀 $x_1$};
 \foreach \r/\c in {1/1,1/2,2/2,2/3,3/3,3/1}{
   \fill[FigureBlue] ({10+\cs*(\c-.5)},{25-\cs*(\r-1)}) circle (1.6);}
 \foreach \r/\c in {1/3,2/1,3/2}{
   \draw[FigureMuted,line width=.6pt,fill=white] ({10+\cs*(\c-.5)},{25-\cs*(\r-1)}) circle (1.2);}
 % table on the right
 \node[figure heading,anchor=west] at (46,35) {可接的末标签};
 \node[figure heading] at (80,35) {$r(u)$};
 \node[figure heading,anchor=west] at (88,35) {计入};
 \foreach \acc/\pair [count=\r] in {{正常，广告}/{$U_{\text{正常},\text{广告}}$},
     {广告，诈骗}/{$U_{\text{广告},\text{诈骗}}$},{正常，诈骗}/{$U_{\text{正常},\text{诈骗}}$}}{
   \draw[FigureMuted,line width=.5pt,-{Stealth[length=1.4mm,width=1mm]}] (41.5,{25-\cs*(\r-1)})--(45,{25-\cs*(\r-1)});
   \node[anchor=west] at (46,{25-\cs*(\r-1)}) {\acc};
   \node at (80,{25-\cs*(\r-1)}) {$2$};
   \node[anchor=west] at (88,{25-\cs*(\r-1)}) {\pair};}
 \node[anchor=west] at (10,-6) {$|\mathcal V|=\textstyle\sum_u r(u)=6
   \;\leq\;|U|+\sum_{\{a,b\}}|U_{a,b}|=3+(1+1+1)$};
\end{tikzpicture}
\endgroup
